\section{Results}
\label{sec:results}

\subsection{Source-only cross-receiver baselines}
Table~\ref{tab:baseline} summarizes source-only performance.
Both directions remain near chance for 24-way file-level classification ($\approx 20\%$).
Source-only prototypes (RCPA-S at $K{=}0$) are similarly weak ($15.3\%\pm 7.3\%$ on RX1$\rightarrow$RX2, $13.0\%\pm 3.1\%$ on RX2$\rightarrow$RX1).

\subsection{RCPA-T main results}
Table~\ref{tab:rcpa} and Fig.~\ref{fig:rcpa_shotcurve} report the RCPA-T shot curve.
Pooled across directions, RCPA-T improves monotonically from $33.1\%\pm 11.2\%$ at $K{=}1$ to $75.0\%\pm 8.0\%$ at $K{=}20$.
At $K{=}5$, pooled accuracy is $58.3\%\pm 9.0\%$, a $+38.2$ percentage-point gain over the source classifier ($\approx 20.1\%$ pooled).
At $K{=}10$, pooled accuracy reaches $69.4\%\pm 9.7\%$.
Direction asymmetry exists but both directions benefit: RX1$\rightarrow$RX2 reaches $57.4\%\pm 9.4\%$ and RX2$\rightarrow$RX1 $59.3\%\pm 8.5\%$ at $K{=}5$.

\begin{table}[t]
\caption{RCPA-T shot curve (file-level accuracy, \%). Mean$\pm$std over 3 seeds$\times$3 splits per direction.}
\label{tab:rcpa}
\centering
\begin{tabular}{ccccc}
\toprule
$K$ & RX1$\rightarrow$RX2 & RX2$\rightarrow$RX1 & Pooled & $\Delta$ vs cls. \\
\midrule
1 & 36.1 ± 11.9 & 30.1 ± 9.4 & 33.1 ± 11.2 & +13.0 pp \\
3 & 49.1 ± 12.5 & 50.5 ± 10.1 & 49.8 ± 11.4 & +29.7 pp \\
5 & 57.4 ± 9.4 & 59.3 ± 8.5 & 58.3 ± 9.0 & +38.2 pp \\
10 & 67.6 ± 10.0 & 71.3 ± 9.1 & 69.4 ± 9.7 & +49.3 pp \\
20 & 77.3 ± 7.1 & 72.7 ± 8.1 & 75.0 ± 8.0 & +54.9 pp \\
\bottomrule
\end{tabular}
\end{table}

\begin{figure}[!t]
  \centering
  \includegraphics[width=\linewidth]{fig2_rcpa_shotcurve}
  \caption{RCPA-T shot curve under block-disjoint $K$-window target-receiver calibration. Points show mean file-level accuracy over three checkpoint seeds and three split repeats; error bars denote standard deviation.}
  \label{fig:rcpa_shotcurve}
\end{figure}

\subsection{Ablations and auxiliary analyses}
Table~\ref{tab:ablation} compares RCPA variants on RX1$\rightarrow$RX2 at $K{=}5$.
RCPA-T ($57.4\%\pm 9.4\%$) strongly outperforms RCPA-S ($15.3\%\pm 7.3\%$) and RCPA-B ($33.8\%\pm 4.1\%$).
Pooled main evaluation at $K{=}10$ shows the same pattern: RCPA-T $69.4\%$ vs.\ RCPA-B $37.0\%$.
Blending source and target prototypes reintroduces source-receiver bias under asymmetric shift.

OOB representation equalization reduces OOB-only receiver probe accuracy from $75.1\%$ to $35.6\%$ with CORAL, confirming that OOB embeddings carry substantial receiver information.
However, this embedding-level alignment does not surpass RCPA-T file-level accuracy in our auxiliary run, supporting its role as mechanism validation rather than a primary method.

\begin{table}[t]
\caption{Ablation and auxiliary analyses (RX1$\rightarrow$RX2 full-mode means at $K$=5 unless noted).}
\label{tab:ablation}
\centering
\begin{tabular}{lc}
\toprule
Method / analysis & File-acc or probe \\
\midrule
RCPA-S (source prototype) & 15.3 ± 7.3 \\
RCPA-T (primary) & 57.4 ± 9.4 \\
RCPA-B (blend ablation) & 33.8 ± 4.1 \\
OOB-Eq CORAL: RX probe before/after & 75.1\% / 35.6\% \\
\bottomrule
\end{tabular}
\end{table}


\subsection{Same-protocol baseline comparison}
Tables~\ref{tab:unlabeled_baselines} and~\ref{tab:kshot_baselines} compare RCPA-T against representative baselines evaluated under the \emph{same} block-disjoint protocol: frozen hybrid embeddings, identical support/query splits, and the same $K$ labeled windows per device.
Direct numerical comparison with prior cross-receiver RFFI systems is not strictly apples-to-apples because those works differ in receiver sets, device counts, source-data availability, target-label assumptions, and calibration protocols~\cite{shen2024receiveragnostic,yang2025scrffi,ma2025receiver_independent_rffi}.
We therefore report same-protocol baselines that cover the adaptation modes most likely to be raised by reviewers.

Unsupervised methods remain far below K-shot calibration.
Feature mean-shift and CORAL alignment on target calibration Block~A, followed by the frozen source classifier, reach only $19.7\%\pm 8.1\%$ and $25.0\%\pm 10.4\%$ pooled file-level accuracy, respectively (Table~\ref{tab:unlabeled_baselines}).
A controlled RX1$\rightarrow$RX2 diagnostic run shows that entropy-min TTA and pseudo-prototype adaptation remain far below RCPA-T (Appendix~\ref{app:tta_diagnostic}, Table~\ref{tab:tta_diagnostic}).
Full source-free adaptation is not the focus of this paper.

Among labeled K-shot methods, RCPA-T is competitive with multinomial logistic regression (linear probe) on frozen embeddings (Table~\ref{tab:kshot_baselines}).
At $K{=}5$, linear probe reaches $59.0\%\pm 12.2\%$ pooled versus $58.3\%\pm 9.0\%$ for RCPA-T---a $0.7$ percentage-point difference within overlapping standard deviations.
At $K{=}10$, RCPA-T reaches $69.4\%\pm 9.7\%$ versus $65.7\%\pm 11.2\%$ for linear probe, with lower cross-split variance.
Classifier-head fine-tuning (source-initialized or random) is consistently below both RCPA-T and linear probe at $K{=}5$ and $K{=}10$ (e.g., $53.9\%\pm 12.0\%$ pooled for source-init head FT at $K{=}5$).
We do \emph{not} claim that RCPA-T comprehensively outperforms all calibration methods; rather, it provides a non-parametric, gradient-free target-receiver calibration option that matches strong K-shot linear baselines while remaining stable at larger $K$.

\begin{table}[t]
\caption{Unlabeled calibration baselines under the same block-disjoint protocol (file-level accuracy, \%). Mean over three checkpoint seeds and three split repeats per direction.}
\label{tab:unlabeled_baselines}
\centering
\small
\begin{tabular}{lccc}
\toprule
Method & RX1$\rightarrow$RX2 & RX2$\rightarrow$RX1 & Pooled \\
\midrule
Source classifier & 19.4$\pm$3.4 & 20.8$\pm$8.8 & 20.1$\pm$6.7 \\
Feat.\ mean-shift + src.\ cls. & 21.8$\pm$4.3 & 17.6$\pm$10.2 & 19.7$\pm$8.1 \\
Feat.\ CORAL + src.\ cls. & 19.4$\pm$5.2 & 30.6$\pm$11.3 & 25.0$\pm$10.4 \\
\bottomrule
\end{tabular}
\end{table}

\begin{table*}[t]
\caption{K-shot labeled calibration baselines under the same block-disjoint protocol (file-level accuracy, \%). Mean over three checkpoint seeds and three split repeats per direction.}
\label{tab:kshot_baselines}
\centering
\small
\begin{tabular}{llcccc}
\toprule
Method & Adapted params & $K$ & RX1$\rightarrow$RX2 & RX2$\rightarrow$RX1 & Pooled \\
\midrule
Linear probe & logistic & 1 & 35.6$\pm$10.0 & 31.0$\pm$10.2 & 33.3$\pm$10.4 \\
Head FT (source init) & linear head & 1 & 33.3$\pm$10.2 & 27.3$\pm$10.0 & 30.3$\pm$10.6 \\
RCPA-T & prototype & 1 & 36.1$\pm$11.9 & 30.1$\pm$9.4 & 33.1$\pm$11.2 \\
Linear probe & logistic & 5 & 57.4$\pm$13.1 & 60.6$\pm$11.0 & 59.0$\pm$12.2 \\
Head FT (source init) & linear head & 5 & 52.8$\pm$13.0 & 55.1$\pm$10.7 & 53.9$\pm$12.0 \\
RCPA-T & prototype & 5 & 57.4$\pm$9.4 & 59.3$\pm$8.5 & 58.3$\pm$9.0 \\
Linear probe & logistic & 10 & 69.4$\pm$13.2 & 62.0$\pm$6.9 & 65.7$\pm$11.2 \\
Head FT (source init) & linear head & 10 & 65.7$\pm$13.0 & 62.0$\pm$6.9 & 63.9$\pm$10.6 \\
RCPA-T & prototype & 10 & 67.6$\pm$10.0 & 71.3$\pm$9.1 & 69.4$\pm$9.7 \\
\bottomrule
\end{tabular}
\end{table*}


\subsection{Diagnosis--result link}
Figure~\ref{fig:diagnosis} connects the diagnosis to the calibration outcome: OOB receiver entanglement and non-transferable source boundaries explain why source-only methods fail, while target-receiver prototypes realign class geometry in the observed receiver space without modifying the backbone.
